
\chapter{Мультивалютность и международные платежи}\label{ch:multicurrency}
Один международный платёж проходит через несколько валют: покупатель видит цену
в~одной, платит картой или~счётом в~другой, продавец получает выплату
в~третьей, а~расчёт с~эквайером приходит в~четвёртой. Каждая сумма фиксируется
в~свой момент времени. Для российского читателя мультивалютность
в~2022--2026~годах означает работу под санкциями: ряд банков отключён от~SWIFT
(Society for Worldwide Interbank Financial Telecommunication~--- международная
сеть финансовых сообщений)~\cite{swift-overview}, корреспондентские отношения
перестраиваются вокруг китайских, белорусских, казахстанских
банков-посредников, растёт доля расчётов в~юанях и~других валютах дружественных
юрисдикций. В~этот контекст входит BRICS (акроним по~странам-учредителям:
  Бразилия, Россия, Индия, Китай, ЮАР; на~2026~год объединение включает
  11~стран, к~учредителям
  присоединились Египет, Эфиопия, Иран, ОАЭ, Саудовская Аравия
и~Индонезия)~\cite{brics-2026-chair}. Корреспондентские счета и~международная
логика карточных сетей разобраны в~предыдущих частях; здесь~--- архитектура
мультивалютного продукта и~валютное регулирование после 2022~года.

\section{Источники курса}\label{sec:mc-currencies}

В~мультивалютном продукте курс фиксирует конкретная сторона; она же отвечает
перед поддержкой и~казначейством.

В~первом варианте курс задаёт эмитент по~правилам карточной сети.
Так устроен стандартный карточный сценарий, когда продавец выставляет цену
в~валюте операции, а~сумма в~валюте карты определяется на~стороне эмитента
по~правилам валюты списания и~конверсии~\cite{visa-core-rules,mc-dcc-guide}.
Для продавца схема проста, но~ни продавец, ни~платформа заранее не~знают
итоговую сумму списания в~валюте карты.

Во~втором варианте курс контролирует эквайер или~провайдер динамической
конвертации валюты (Dynamic Currency Conversion, DCC), который показывает
держателю карты итоговую сумму в~домашней валюте ещё до~авторизации. Цена
становится предсказуемой для покупателя, но~право задать курс и~удержать маржу
переходит к~стороне, которая управляет DCC~\cite{visa-core-rules,mc-dcc-guide}.

В~третьем варианте курс задаёт сама платформа. Так устроен кошелёк, сервис
путешествий или~маркетплейс с~собственным валютным модулем, где интерфейс
показывает курс, платформа берёт на~себя риск его устаревания за~короткий
интервал, а~затем хеджирует (страхует валютный риск встречным инструментом)
или~неттирует (взаимозачёт встречных требований до~фактического перевода~---
\enquote{неттинг}) фактические потоки на~этапе расчёта.

В~четвёртом варианте курс предоставляет внешний \textbf{FX-провайдер} (foreign
  exchange~--- операции на~валютном рынке; источник зависит от~валютной пары
и~канала~--- биржевая площадка, межбанковский рынок или~уполномоченный банк),
а~платёжный продукт привязывает платёж к~уже рассчитанной сделке. Сценарий
типичен для B2B-сервисов (business-to-business, расчёты между юридическими
лицами) и~платформ выплат, где цена перевода действует дольше одной сессии
оплаты. Валютный риск остаётся на~казначейской функции корпоративного клиента;
отдельная сделка, например форвардный контракт, закрывает этот риск.

Платформа показывает клиенту курс в~момент выбора, подтверждение приходит через
час, и~за~это время курс изменился. Разница между показанным и~фактическим
курсом~--- убыток той стороны, которая зафиксировала обязательство. Курс
действует в~пределах TTL (time-to-live~--- срок действия котировки),
и~при~позднем подтверждении платформа либо~пересчитывает по~актуальному курсу,
либо~принимает риск на~себя. Значения TTL и~допустимый диапазон волатильности
для основных валютных пар оператор подбирает по~собственным данным; короткий
TTL увеличивает долю истёкших сессий, длинный~--- валютный риск. Без~явного TTL
потери на~устаревшем курсе обнаруживаются только при~ежемесячной сверке
казначейства (treasury~--- внутренней функции компании, отвечающей
  за~управление денежной позицией, ликвидностью и~валютным риском; в~отличие
от~государственного казначейства, упомянутого в~главе~\ref{ch:cbdc}).

\begin{figure}[htbp]
  \centering
  \includefiguresvg[width=.65\linewidth]{assets/figures/ch35-multicurrency/fx-ownership-models}
  \caption{Четыре модели владения курсом в~мультивалютном
  продукте.}\label{fig:fx-ownership-models}
\end{figure}

Подписка с~ежемесячным списанием в~EUR при~выплате продавцу в~USD проходит
через конверсию 12~раз в~год, и~при~каждой конверсии курс другой; совокупная
дельта влияет на~маржу платформы. Курс хранится как отдельный объект с~явным
источником, значением, моментом фиксации, TTL и~стороной, несущей риск.

Минимальный реестр котировок принимает котировки от~разных источников,
отбрасывает устаревшие (по~TTL) и~слишком широкие по~спреду (выше заданного
  порога в~базисных пунктах, basis points, bps~--- сотых долях процентного
пункта), и~возвращает последнюю по~времени из~оставшихся. Источник времени
вынесен в~параметр \texttt{now}, чтобы тестируемое решение не~зависело
от~системного таймера:

\practicenote{%
  Если продукт обещает клиенту \enquote{ваш курс}, у~этого курса должны
  быть явный источник и~срок действия.

}

\needspace{4\baselineskip}
\codefile{Python}{samples/ch35-fx-quote/quote.py}

\FloatBarrier
\section{Фиксация валюты расчёта}\label{sec:mc-settlement}

\textbf{Валюта цены} (transaction currency, в~правилах карточных схем также
\enquote{валюта операции})~--- та, которую видит покупатель и~которую продавец
ожидает на~странице оплаты. \textbf{Валюта списания} (billing currency, она же
\enquote{валюта карты})~--- та, в~которой эмитент дебетует держателя карты.
\textbf{Валюта расчёта} (settlement currency)~--- та, в~которой сеть
или~эквайер фактически рассчитывается с~вашей
стороной~\cite{visa-core-rules,mc-rules}. Трёхбуквенные коды валют (RUB, USD,
EUR, CNY, KZT и~др.) даются по~стандарту ISO~4217~\cite{iso-4217}.

Держатель карты работает в~своей валюте, расчётная инфраструктура схемы~---
в~валюте расчёта, эквайер~--- в~валюте своей юрисдикции. Для международных
коридоров в~Visa и~Mastercard валютой расчёта по~умолчанию выступают USD
или~EUR; в~платёжной системе \enquote{Мир} внутрироссийский расчёт проходит
в~рублях через платёжную систему Банка России.

При \textbf{единой валюте расчёта} внутренний реестр обязательств и~выплаты
ведутся в~одной базовой валюте, поэтому учёт упрощается до~одного казначейского
баланса, одной логики комиссий и~меньшего числа сценариев для~сверки. Платой
за~это становится постоянная конверсия на~входе или~выходе из~системы
и~зависимость от~валютной наценки внешнего провайдера.

При \textbf{нескольких валютах расчёта} система хранит обязательства и~расчёты
в~нескольких валютах. Реестр обязательств и~логика выплат сложнее, но~лишней
конверсии нет, а~фактические валютные позиции не~смешиваются в~одну
синтетическую сумму.

Один и~тот же EUR в~разных частях системы~--- это валюта цены, валюта списания,
валюта расчёта или~валюта хранения обязательства; роли могут не~совпадать.

Продавец из~России продаёт цифровой продукт покупателю из~Казахстана; цена
выставлена в~рублях, покупатель оплачивает \textbf{картой UnionPay
казахстанского банка} в~тенге, эквайер рассчитывается с~продавцом в~рублях
через российского эквайера UnionPay. Валюта цены~--- RUB, валюта списания~---
KZT (валюта карты), валюта расчёта~--- RUB. После марта 2022~года это
\textit{единственный массовый карточный канал} для иностранного держателя карты
в~российский ТСП (торгово-сервисное предприятие): международные сети Visa
и~Mastercard в~РФ закрыты для иностранных карт~\cite{cbr-ps,nspk-mir-rules},
а~мосты НСПК с~локальными схемами маршрутизируют только пары
\enquote{Мир~$\leftrightarrow$~локальная~ПС}, не~карты Visa или~Mastercard. Из~четырёх
таких мостов на~дату издания устойчиво работает лишь БЕЛКАРТ (Беларусь):
Uzcard (Узбекистан) прекратил приём \enquote{Мира} ещё в~2022~году, а~АРКА
(Армения) и~Элкарт (Кыргызстан) свернули его в~2024~году вслед за~SDN-листингом
НСПК.
Реестр обязательств должен хранить все точки фиксации~--- сумму в~KZT на~момент
авторизации, применённый курс при~клиринге и~сумму в~RUB на~момент расчёта.

У~каждой фазы свой источник курса: \mbox{EUR$\to$USD} конвертируется правилами
схемы при~клиринге, \mbox{USD$\to$GBP}~--- внешним FX-провайдером выплаты
на~момент перевода продавцу. Авторизация и~клиринг занимают секунды и~часы,
расчёт и~выплата~--- дни.

\begin{figure}[htb]
  \centering
  \includefiguresvg[width=.8\linewidth]{assets/figures/ch35-multicurrency/three-currencies-timeline}
  \caption{Жизненный цикл операции в~нескольких валютах: transaction
  $\to$ billing $\to$ settlement $\to$ payout.}\label{fig:three-currencies-timeline}
\end{figure}

\section{DCC и~маршрутизация по~валютным коридорам}\label{sec:mc-dcc}

\textbf{Валютный коридор}~--- пара \enquote{страна покупателя / страна ТСП}
или~\enquote{валюта источника / валюта получателя}, для которой задан набор
маршрутов, эквайеров и~провайдеров. Коридор определяет, через какую сеть
пойдёт авторизация и~в~какой валюте пройдёт расчёт.

DCC показывает держателю карты сумму в~знакомой валюте. Расчёт по~сети идёт
по~маршруту схемы, а~сумма списания для клиента считается
отдельно~\cite{visa-core-rules,mc-tpr}.

Продукт хранит исходную сумму, показанную сумму и~применённый курс как
отдельные значения, не~выводя одно из~другого задним числом. Без~этого нельзя
объяснить клиенту списание и~сверить удержанную маржу
DCC~\cite{visa-core-rules,mc-tpr}.

Маржа DCC возникает на~спреде (разнице между ценой покупки и~ценой продажи
валюты), когда конверсия предлагается по~курсу, менее выгодному для держателя
карты, чем межбанковский (он~же оптовый~--- курс, по~которому платёжная сеть
конвертирует на~стороне расчёта). Разница в~виде наценки (валютной маржи сверх
оптового курса) остаётся у~эквайера или~его провайдера DCC. Конкретное
распределение определяется коммерческим договором между ними и~в~публичных
правилах схем не~закреплено. Держатель карты видит сумму в~знакомой валюте
до~операции, но~платит за~это наценкой, встроенной в~курс. Схемы требуют явного
раскрытия наценки и~согласия клиента до~операции.

Покупатель платит 100~EUR картой европейского банка в~ТСП в~России. Без~DCC
эквайер отправляет транзакцию на~100~EUR; сеть конвертирует её при~клиринге
по~своему оптовому курсу с~наценкой сети, и~держатель карты видит итоговую
сумму в~своей домашней валюте только после списания. С~DCC эквайер через
провайдера DCC предлагает держателю карты фиксированную сумму в~его домашней
валюте сразу на~терминале по~курсу провайдера; разница между этой наценкой
и~оптовым курсом~--- плата за~фиксацию суммы до~операции. Значения наценок
отличаются у~разных провайдеров DCC и~публикуются в~раскрытии перед операцией.

Если ТСП продаёт только в~одной стране и~получает выплаты в~местной валюте, DCC
не~нужен: конверсию делает сеть при~клиринге. Если ТСП
хочет предложить DCC ради повышения маржи, нужен провайдер DCC, раскрытие
наценки и~согласие клиента; репутационный риск у~этого варианта высокий. Если
ТСП работает в~нескольких странах и~хочет снизить долю международных отказов,
локальный эквайер в~каждой стране повышает долю одобрений и~убирает
международную наценку, но~требует юридического присутствия и~отдельной
расчётной модели в~каждой юрисдикции~\cite{visa-core-rules,mc-tpr}.

Политика маршрутизации по~коридору, BIN, стране эквайера и~типу~ТСП задаёт
допустимость DCC и~источник курса.

Согласие клиента, раскрытая наценка и~данные для последующей проверки
фиксируются как поля события транзакции, иначе позже трудно доказать,
что~конверсия была предложена корректно~\cite{visa-core-rules,mc-tpr}.

Платформа, работающая с~несколькими странами, эквайерами и~схемами выплат,
маршрутизирует платежи по~коридорам: для~каждой пары \enquote{страна покупателя
/ страна ТСП} выбираются эквайер, валютная схема и~способ выплаты.

Выбор коридора зависит от~стоимости авторизации и~расчёта, доли одобрений для
нужного диапазона BIN, валюты выплаты продавцу и~стороны, которая несёт
валютный риск. Параметры обоих маршрутов по~этим факторам показаны
на~рис.~\ref{fig:corridor-routing}; состав параметров согласован с~публичными
правилами Visa и~таблицами Visa USA
IRF~\cite{visa-core-rules,visa-interchange}.

Дорогой маршрут через эквайера страны держателя карты повышает
\(P(\text{approve})\), но~увеличивает комиссию. Дешёвый маршрут через
эквайера страны ТСП снижает комиссию, но~и~\(P(\text{approve})\)
на~международных картах становится ниже. При~высоком среднем чеке потеря
на~отказанной транзакции дороже разницы в~комиссии, и~выбор смещается
к~дорогому маршруту; при~низком чеке выгоднее обратное.

При \textbf{выплате в~домашней валюте продавца} платформа сама консолидирует
и~конвертирует все входящие потоки. Для продавца модель проще, для
казначейской функции платформы~--- сложнее.

При \textbf{выплате в~валюте исходного потока} платформа ведёт
\textbf{раздельные валютные балансы продавца} (отдельный баланс на~каждую
валюту) и~платит продавцу в~той валюте, в~которой пришёл поток. Конверсий
меньше, но~у~продавца появляется несколько валютных балансов и~отдельный
график выплат по~каждому.

\begin{figure}[htbp]
  \centering
  \includefiguresvg[width=.85\linewidth]{assets/figures/ch35-multicurrency/corridor-routing}
  \caption{Маршрутизация по~коридору: сравнение двух стратегий
  для одного платежа.}\label{fig:corridor-routing}
\end{figure}

\practicenote{%
  Если внутри принимаются потоки в~нескольких валютах, а~продавцу обещан один
  баланс, у~валютной позиции (чистой величины обязательств в~данной
  валюте~--- длинной или~короткой), лимитов и~источников курса должен
  быть явный владелец.

}

Свести выбор коридора к~одной метрике нельзя. Доля одобрений, комиссия, валюта
выплаты и~сторона валютного риска~--- четыре независимых параметра; решение,
лучшее по~одному из~них, редко остаётся лучшим по~остальным. Платформа
закрепляет за~каждой парой \enquote{страна покупателя / страна ТСП} одну
рабочую конфигурацию и~поддерживает резервные: при~падении доли одобрений
или~при~отзыве BIN-диапазона переход на~резервный маршрут должен занимать
минуты, а~не недели. Платформа считает фактическую стоимость каждой
конфигурации по~среднему чеку: на~низком чеке экономия комиссии перекрывает
потерю на~отказах, на~высоком~--- наоборот. Цифры берутся из~собственной
статистики, потому что публичные таблицы interchange покрывают только
комиссионную часть; поведение конкретного BIN на~конкретном эквайере наблюдает
только оператор.

\FloatBarrier
\section{Сверка по~валютам}\label{sec:mc-fx-risk}

Расхождения в~мультивалютном продукте появляются после авторизации, когда одна
и~та же операция представлена сразу в~нескольких представлениях: на~странице
оплаты~--- в~валюте цены, в~журнале авторизации~--- одной суммой, в~клиринге~---
с~другой валютой расчёта, в~отчёте о~выплате~--- после удержания комиссии
и~валютного спреда. Валютная сверка остаётся отдельным процессом: карточные
правила прямо требуют сверять данные сети с~внутренними записями,
а~в~Mastercard это сформулировано как ежедневная обязанность сверять итоги
и~транзакции с~внутренним учётом~\cite{mc-rules}.

\highlight{Внутренний реестр обязательств хранит исходную валютную пару
  полностью~--- исходную сумму, сумму после конверсии, источник курса и~момент
его фиксации}.

Комиссия и~конверсия считаются отдельно. Если курсовая разница смешана
с~комиссией эквайера, казначейство не~увидит, какая часть маржи ушла в~курс,
а~какая~--- в~комиссию.

Выплаты сверяются по~раздельным валютным балансам, а~не по~агрегированной
сумме. Совпавший общий итог скрывает недостачу в~отдельной валюте, когда одна
валюта компенсирует другую.

Покупка за~100~EUR: расчёт по~курсу 1,08 = 108~USD, комиссия эквайера 2,5\,\% =
2,70~USD, interchange 1,5\,\% = 1,62~USD, комиссия платформы 3\,\% = 3,24~USD,
выплата продавцу = 108 $-$ 2,70 $-$ 1,62 $-$ 3,24 = 100,44~USD. Если продавец
получает выплату в~GBP по~курсу 0,79, итого 79,35~GBP. Одна и~та же покупка
существует в~пяти представлениях~--- 100~EUR на~странице оплаты, 100~EUR
в~авторизации, 108~USD в~расчёте, 100,44~USD в~нетто-выплате, 79,35~GBP
на~счёте продавца.

Типовая ошибка сверки~--- смешение валютной дельты и~комиссий. Если 7,70~USD
удержаний (2,70 + 1,62 + 3,24) учтены как курсовой убыток вместо раздельных
комиссий, казначейство получает завышенную курсовую дельту. Каждое удержание
имеет свой тип: interchange, assessment, комиссия платформы, валютная
наценка~--- и~сверяется отдельно; агрегированная сумма скрывает эту структуру.
На~рис.~\ref{fig:multicurrency-recon} три удержания свёрнуты в~единую дельту
между расчётом и~нетто-выплатой; валюта списания держателя карты (billing)
на~схеме отдельной точкой не~показана.

\vspace{-2pt}
\begin{figure}[htbp]
  \centering
  \includefiguresvg[width=.85\linewidth]{assets/figures/ch35-multicurrency/multicurrency-recon}
  \caption{Одна транзакция в~пяти валютных
  представлениях.}\label{fig:multicurrency-recon}
\end{figure}

\section{Валютное регулирование после 2022~года}\label{sec:mc-russia}

Валютные операции в~России регулирует Федеральный закон от~10.12.2003~№~173-ФЗ
\enquote{О~валютном регулировании и~валютном контроле}~\cite{fz-173}.
Он~определяет понятия резидента и~нерезидента, перечень разрешённых валютных
операций между ними и~порядок валютного контроля. Подзаконные акты Банка России
конкретизируют три процедуры. \textbf{Представление документов валютного
контроля}~--- передача в~уполномоченный банк контракта, счетов, актов и~прочих
подтверждающих документов по~операциям с~нерезидентами. \textbf{Постановка
контракта на~учёт}~--- присвоение контракту уникального номера в~упол\-номоченном
банке; обязательна для контрактов с~нерезидентами выше пороговой суммы,
закреплённой Инструкцией Банка России от~16.08.2017~№~181-И
\enquote{О~представлении документов валютного
контроля}~\cite{cbr-181i}.\footnote{Инструкция 181-И заменила Инструкцию
138-И.}
\textbf{Репатриация валютной выручки}~--- обязанность резидента-экспортёра
обеспечить поступление валютной выручки на~счета в~уполномоченных банках,
а~резидента-импортёра~--- возврат авансов по~неисполненным импортным контрактам
в~срок, установленный контрактом. В~платёжном продукте с~международными
операциями с~участием российской стороны 173-ФЗ и~подзаконные акты учитываются
при~открытии валютных счетов, идентификации клиента, направлении валютной
операции в~банк и~формировании отчётности.

После марта 2022~года международные валютные расчёты для российских банков
перестроились: ряд крупнейших банков был отключён от~SWIFT по~решению
Совета~ЕС, оформленному Регламентом Совета (EU)
2022/345~\cite{eu-swift-russia-2022,ofac-russia-sanctions}. SWIFT работает для
значительной части российских банков и~для большинства корреспондентских
отношений с~банками дружественных стран, но~отключённым банкам пришлось
перенести финансовые сообщения на~альтернативные каналы.

\textbf{СПФС} (Система передачи финансовых сообщений Банка России;
см.~раздел~\ref{sec:banks-infrastructure}) для отключённых от~SWIFT банков
превратилась из~локальной альтернативы в~основной канал финансовых сообщений
с~банками-партнёрами в~ЕАЭС, СНГ и~BRICS~\cite{cbr-spfs}.

\textbf{Расчёты в~юанях} стали отдельным каналом для экспорта и~импорта после
2022~года. Расчёты в~CNY между российскими и~китайскими банками проходят через
систему CIPS (Cross-Border Interbank Payment System)~---
китайскую инфраструктуру межбанковских расчётов в~юанях. Оператор~---
Cross-border Interbank Payment System Co., Ltd., под~надзором Народного банка
Китая~\cite{cips-overview}. Российские банки подключаются к~CIPS как прямые
или~непрямые участники либо проводят платежи через банки-посредники в~Китае.

\textbf{Корсчета через банки-посредники.} Для отключённых от~SWIFT российских
банков корреспондентские отношения с~зарубежными банками-партнёрами
перестраиваются через банки-посредники в~дружественных юрисдикциях~--- Китай,
Индия, ОАЭ, Турция, страны Центральной Азии. Цепочка расчётов удлиняется,
каждая операция проходит дополнительные проверки по~комплаенсу и~валютному
контролю.

\textbf{BRICS Pay и~связанные инициативы.} В~BRICS обсуждаются проекты
межоператорной интеграции национальных платёжных систем (BRICS Pay, BRICS
Bridge) для прямых расчётов в~национальных валютах без~обращения
к~доллару~\cite{brics-pay-overview}.\footnote{Отдельный проект межбанковского
  обмена цифровыми валютами центральных банков (CBDC, central bank digital
  currency) mBridge~--- многосторонняя платформа CBDC-расчётов~--- вёлся под~BIS
  Innovation Hub (Банк международных расчётов, Bank for International
  Settlements) совместно с~центральными банками Китая, ОАЭ, Гонконга и~Таиланда;
  BIS вышел из~проекта 31.10.2024, дальнейшее развитие продолжают участники
без~BIS; mBridge не~является проектом BRICS.}

Для архитектуры платёжного продукта в~России это даёт три последствия.
Во-первых, мультивалютная модель должна предусматривать несколько
каналов расчёта: SWIFT (для неотключённых банков и~зарубежных коридоров),
СПФС (для внутрироссийских расчётов и~ЕАЭС), CIPS (для расчётов в~юанях),
двусторонние и~многосторонние соглашения. Во-вторых, после 2022~года выбор
валют расчёта сместился к~национальным валютам дружественных стран (CNY, AED,
INR, KZT, BYN); USD и~EUR остаются возможными, но~требуют согласования
с~банком-партнёром по~каждой
операции и~подпадают под~санкционный контроль. В-третьих, валютный контроль
по~173-ФЗ требует представления документов и~постановки контрактов на~учёт
независимо от~выбранного канала расчёта.

\practicenote{%
  Обещание \enquote{международный платёж в~любой валюте} должно соответствовать
  фактическим валютным коридорам: набор работающих валют, банков-посредников
  и~каналов зависит от~санкционного статуса участников и~меняется.

}

\section{Корреспондентский банкинг и ISO~20022}\label{sec:mc-correspondent}

Корреспондентские расчёты разобраны в~разделе~\ref{sec:banks-correspondent},
миграция к~ISO~20022~--- в~разделе~\ref{sec:iso20022}.

Расчёт по~международной операции проходит по~более длинной цепочке, чем
авторизация, и~занимает больше
времени~\cite{fatf-correspondent-banking,swift-instant-interoperability}:
одобрение за~секунду не~означает выплату в~тот~же день. Расчётный банк,
карточная сеть и~валютный провайдер в~модели продукта фиксируются как разные
роли; продукт указывает ответственного за~каждую стадию без~отговорок вида
\enquote{деньги идут через~SWIFT}.

Переход отрасли на~ISO~20022 даёт структурированные поля, которые облегчают
комплаенс-проверки и~сверку, но~сами по~себе не~отвечают, кто несёт валютный
риск и~в~какой валюте учитываются внутренние
обязательства~\cite{iso-20022}.\footnote{\cite{swift-iso20022-milestone,swift-transaction-manager-iso20022}.}

Структурированные поля ISO~20022 дают продукту явное представление участников
перевода~--- но~только когда внутренняя модель уже различает плательщика,
конечного получателя, посредников и~назначение перевода как отдельные сущности.
Если в~исходной модели платёж сводился к~паре \enquote{отправитель~---
получатель} с~произвольным текстом в~поле описания, миграция к~ISO~20022
на~стороне сетевых сообщений не~решит проблему: внутри продукт по-прежнему
теряет цепочку. Поэтому переход начинается с~инвентаризации внутренних полей,
а~не с~шлюзов.

Корреспондентский банкинг в~санкционных условиях создаёт ещё одно
требование: каждое дополнительное звено цепочки (банк-посредник в~дружественной
юрисдикции) добавляет комплаенс-проверки и~риск возврата платежа на~любой
стадии. Продукт фиксирует, кто и~какую информацию отправляет на~каждом звене,
и~хранит копию переданных полей на~своей стороне~--- иначе для разбора отказа
придётся восстанавливать цепочку сообщений по~логам трёх-четырёх банков
с~разными политиками хранения. ISO~20022 формализует синтаксис, но~не отменяет
необходимости вести собственный реестр стадий расчёта.
